\section{Finite tori and separated minimizing configurations}
\label{geo:section}\label{sec:geometry}

We first isolate the finite-dimensional problem used in the periodic
reduction.  Separation will be proved only for its minimizing
configurations.  This suffices for a lower bound on every configuration.

For an integer \(n\geq2\), put
\[
 L=\sqrt n,\qquad T_n=\R^2/(L\Z^2).
\]
The torus has area \(n\).  Given distinct points
\(v_1,\ldots,v_n\in T_n\), let \(h\) be the mean-zero distributional
solution of
\[
 -\Delta h=2\pi\left(\sum_{i=1}^n\delta_{v_i}-1\right).
\]
We use the energy \eqref{eq:torus-energy}. For this square-torus
energy, \(\eta\) is smaller than the injectivity radius and half the
minimum distance between the finitely many points.

\subsection{Green representation and existence}

\begin{lemma}[Green representation and existence]
\label{geo:existence}
Let \(G_n\) be the mean-zero Green function on \(T_n\), normalized by
\[
 -\Delta G_n=2\pi(\delta_0-1/n),
 \qquad
 R_n=\lim_{x\to0}\bigl(G_n(x)+\log|x|\bigr).
\]
For every configuration of \(n\) distinct points,
\begin{equation}
 h(x)=\sum_{i=1}^nG_n(x-v_i),\qquad
 \mathcal W_n(h)=\pi nR_n+
                    \pi\sum_{i\ne j}G_n(v_i-v_j).
 \label{geo:pair-energy}
\end{equation}
The minimum of \(\mathcal W_n\) over such configurations is attained
at a configuration of distinct points.
\end{lemma}

\begin{proof}
The first identity follows from the equation and the zero-mean
normalization.  Subtracting the logarithmic singularity and applying
local elliptic regularity gives, near \(v_i\),
\[
 h(x)=-\log|x-v_i|+a_i+O(|x-v_i|),\qquad
 a_i=R_n+\sum_{j\ne i}G_n(v_i-v_j).
\]
Let \(\Omega_\eta=T_n\setminus\bigcup_iB(v_i,\eta)\).
On its hole boundaries the outward normal points towards \(v_i\);
therefore
\[
 \partial_{\boldsymbol n}h=\frac1\eta+O(1).
\]
Since \(\Delta h=2\pi\) on \(\Omega_\eta\), integration by parts gives
\[
 \frac12\int_{\Omega_\eta}|\nabla h|^2
 =-\pi\int_{\Omega_\eta}h
   +\frac12\sum_i\int_{\partial B(v_i,\eta)}
                         h\,\partial_{\boldsymbol n}h.
\]
The bulk integral tends to zero because \(h\) has mean zero.  The
\(i\)-th boundary integral is
\(2\pi(-\log\eta+a_i)+o(1)\).  Taking the puncture limit proves
\eqref{geo:pair-energy}.

The function \(G_n\) is bounded below on its punctured torus and tends
to \(+\infty\) at the origin.  Extend the right-hand side of
\eqref{geo:pair-energy} to \(+\infty\) on configurations with a
collision.  This gives a lower semicontinuous function on the compact
space \(T_n^n\), finite at every configuration of distinct points.
A collision cannot be canceled by other pair terms, since all those
terms have a common finite lower bound.  A minimum therefore exists
and has distinct points.
\end{proof}

By \eqref{red:harmonic}, every admissible periodic field with these
charges differs from \(-\nabla h\) by a constant vector \(c\), which
adds \(n|c|^2/2\) to the energy.  It therefore suffices to minimize
over the mean-zero potentials considered here.

\subsection{Separation of minimizers}

The argument uses a screened logarithmic potential and the minimum
principle. A closely related exclusion argument for minimizing
configurations of a planar confined system, with the same single-charge
screening radius, appears in Lieb, Rougerie, and
Yngvason~\cite[Lemma~3 and Appendix~A]{LiebRougerieYngvason2018}.
The following proof treats the torus Green function and its background
directly.

\begin{lemma}[Separation]
\label{geo:separation}
For every \(n\geq2\), a minimizing configuration in
Lemma~\ref{geo:existence} has torus distance at least
\[
 r_*=\frac1{\sqrt\pi}
\]
between distinct points.  Its full lift to \(\R^2\) has the same
minimum-separation bound, including between different period
translates of one torus point.
\end{lemma}

\begin{proof}
Fix one point \(q\) of a minimizing configuration.  The symmetry of
\(G_n\), together with \eqref{geo:pair-energy}, shows that
\[
 \Phi(x)=\sum_{p\ne q}G_n(x-p)
\]
has a global minimum at \(q\).  At a fixed point \(p\ne q\), interpret
\(\Phi(p)=+\infty\).

The injectivity radius of \(T_n\) is \(\sqrt n/2>r_*\) for \(n\geq2\).
We may therefore define a function on the embedded radius-\(r_*\)
disk by
\[
 \phi(x)=
 \begin{cases}
 -\log(|x|/r_*)+\dfrac{\pi}{2}(|x|^2-r_*^2),
       &0<|x|<r_*,\\[3pt]
 0,&|x|\geq r_*.
 \end{cases}
\]
Its value and radial derivative vanish at \(r_*\), and it is
positive inside the disk.  Consequently
\[
 -\Delta\phi=2\pi(\delta_0-\mathbf1_{B(0,r_*)})
\]
without a boundary measure on the circle.  Each such disk has area
one.

Set
\[
 \Omega=\bigcup_{p\ne q}B(p,r_*),\qquad
 \Psi=\Phi-\sum_{p\ne q}\phi(\,\cdot-p).
\]
All poles cancel.  In particular \(\Psi\in W^{2,s}(T_n)\) for every
finite \(s\), and \(\Psi\) is continuous.  On \(\Omega\),
\[
 -\Delta\Psi
 =2\pi\sum_{p\ne q}\mathbf1_{B(p,r_*)}
      -2\pi\frac{n-1}{n}
 \geq \frac{2\pi}{n}>0.
\]
Since \(|\Omega|\leq n-1<n\), every component of \(\Omega\) has
nonempty boundary.  Every obstacle vanishes on this boundary, so
\(\Psi=\Phi\) there.

If \(q\in\Omega\), let \(D\) be its component.  The strict minimum
principle, valid for this continuous \(W^{2,s}\) supersolution,
gives
\[
 \Phi(q)>\Psi(q)>
       \min_{\partial D}\Psi
       =\min_{\partial D}\Phi\geq\Phi(q),
\]
a contradiction.  No regularity of the union's boundary is required:
the minimum is attained on \(\overline D\), and strict
superharmonicity excludes an interior minimum.  Hence
\(q\notin\Omega\).  As \(q\) was arbitrary, distinct torus points
are separated by \(r_*\).

Distances between lifts of distinct torus points are at least
their torus distance.  Two different lifts of one point differ by
a nonzero vector in \(L\Z^2\), whose length is at least
\(L=\sqrt n>r_*\).
\end{proof}

\section{Voronoi sectors and a glued dual potential}
\label{dual:section}

We first record the sector geometry of separated square-torus
minimizers. We then prove gluing and duality for compatible periodic
Voronoi sectors on any flat torus. This separates the general identity
from its square-torus and triangular applications. The construction
uses general radial functions; their particular choice and triangular
calibration will be given later. Its purpose is to turn a sector lower
bound, established by two scalar estimates, into an energy bound for a
whole configuration.

The use of polygon inequalities together with an average of at most
six sides is familiar from moment methods; see Gruber's proof of
Fejes T\'oth's theorem~\cite{Gruber1999}. A related affine cell bound
for a Wasserstein interaction is sharp at the regular hexagon of unit
mass; the result of Bourne, Peletier, and
Theil~\cite[arXiv v3, Lemma~12]{BournePeletierTheil2014} assumes integer
side counts at least three and cell mass \(v\geq1.5\cdot10^{-4}\).
Here we must construct a globally admissible singular potential before
estimating its sector contributions. The arguments below supply that
construction and the incidence count on a torus.

\subsection{Counting sectors on the torus}

For \(d>0\) and real \(\alpha,\beta\) with
\(|\alpha|,|\beta|<\pi/2\) and \(\alpha+\beta>0\), define the sector
\begin{equation}
 S(d,\alpha,\beta)=
 \{(r,\varphi):-\alpha<\varphi<\beta,\ 
                      0<r<d\sec\varphi\}.
 \label{geo:sector}
\end{equation}
Its polar axis is the normal from its site toward the supporting edge.

\begin{lemma}[Periodic Voronoi geometry]
\label{geo:voronoi}
Let \(P\subset\R^2\) be the full \(L\Z^2\)-periodic lift of \(n\)
distinct points on \(T_n\), and suppose \(P\) has separation at
least \(r_*\).  Its Voronoi polygons, modulo \(L\Z^2\), have at most
\(6n\) face--edge incidences.

Each such incidence determines a sector \eqref{geo:sector} in polar
coordinates about its site, with \(d\geq r_*/2\).
The sectors in one torus have total area \(n\) and total angle
\(2\pi n\).  Either endpoint parameter \(\alpha\) or \(\beta\) is
allowed to be negative.
\end{lemma}

\begin{proof}
For \(p\in P\), the Voronoi cell
\[
 C_p=\{x:|x-p|\leq|x-q|\text{ for all }q\in P\}
\]
contains the disk of radius \(r_*/2\) about \(p\).  Comparison with
the four sites \(p\pm(L,0)\) and \(p\pm(0,L)\) gives
\[
 C_p\subset p+[-L/2,L/2]^2.
\]
Local finiteness now implies that \(C_p\) is a bounded convex
polygon with finitely many sides.  There are finitely many
polygon, edge, and vertex orbits under the period lattice.

Use maximal positive-length Voronoi edges.  Each edge has two
face incidences, and each vertex has degree at least three.
Indeed a Voronoi vertex has at least three nearest sites; a
degenerate vertex with more nearest sites only increases the
degree.  The quotient is a finite cell decomposition of the
torus with \(F=n\) open faces.  Each open face is a disk: distinct
translates of the interior of a lift cell are disjoint.
A projected closed face need not be embedded, since its boundary
may identify to itself.  These identifications are permitted in
the cell decomposition.

Let \(V\) and \(E\) denote its numbers of vertices and edges.
Euler's identity and the degree count give
\[
 V-E+n=0,\qquad 3V\leq2E.
\]
Thus \(E\leq3n\), and the number of face--edge incidences is
\(M_0=2E\leq6n\).  These counts include every occurrence in a
face boundary: an edge may be a loop, a face may lie on both
sides of an edge, a neighbor may recur, and a vertex may occur
several times in one boundary.  No generic-position assumption
is used.

For an edge of \(C_p\) shared with \(C_q\), choose as polar axis
the unit vector from \(p\) towards \(q\).  Its supporting line
has distance \(d=|p-q|/2\geq r_*/2\) from \(p\).
The angular coordinates of its two endpoints satisfy
\(-\pi/2<\varphi_-<\varphi_+<\pi/2\).  Write
\(\varphi_-=-\alpha\) and \(\varphi_+=\beta\).
The triangle from \(p\) to the edge is exactly
\eqref{geo:sector}.  If the perpendicular foot is outside the
edge, one of \(\alpha,\beta\) is negative.

The sectors partition each polygon up to their boundary rays.
A polygon has angle \(2\pi\) about its interior site.  Summing
areas and angles over the \(n\) quotient faces proves the last
assertions.
\end{proof}

For reference, the area and angle of \eqref{geo:sector} are
\begin{equation}
 A(S)=\frac12\int_{-\alpha}^{\beta}d^2\sec^2\varphi\,\dd\varphi,
 \qquad
 \theta(S)=\alpha+\beta.
 \label{geo:area-angle}
\end{equation}
In particular the sector has positive angle even when one endpoint
parameter is negative.

\begin{figure}[tb]
\centering
\begin{tikzpicture}[scale=.95,>=Stealth,font=\small]
\begin{scope}[shift={(0,0)}]
  \fill[blue!9] (0,0)--(2.4,-1.05)--(2.4,1.55)--cycle;
  \draw[thick] (0,0)--(2.4,-1.05)--(2.4,1.55)--cycle;
  \draw[dashed,gray] (0,0)--(2.9,0);
  \draw[gray] (2.2,0)--(2.2,.2)--(2.4,.2);
  \fill (0,0) circle (1.5pt) node[left] {$p$};
  \node[right] at (2.4,-1.05) {$v_-$};
  \node[right] at (2.4,1.55) {$v_+$};
  \draw[<->] (0,-1.65)--node[above] {$d$}(2.4,-1.65);
  \draw[->] (.72,0) arc[start angle=0,end angle=32.86,radius=.72];
  \node at (.96,.30) {$\beta$};
  \draw[->] (.64,0) arc[start angle=0,end angle=-23.63,radius=.64];
  \node at (.94,-.27) {$-\alpha$};
  \node[align=center] at (1.15,-2.35) {The foot lies on the edge.\\$\alpha,\beta\geq0$};
\end{scope}
\begin{scope}[shift={(6.3,0)}]
  \fill[blue!9] (0,0)--(2.4,.65)--(2.4,2)--cycle;
  \draw[thick] (0,0)--(2.4,.65)--(2.4,2)--cycle;
  \draw[dashed,gray] (0,0)--(2.9,0);
  \draw[dashed,gray] (2.4,0)--(2.4,.65);
  \draw[gray] (2.2,0)--(2.2,.2)--(2.4,.2);
  \fill (0,0) circle (1.5pt) node[left] {$p$};
  \node[right] at (2.4,.65) {$v_-$};
  \node[right] at (2.4,2) {$v_+$};
  \draw[<->] (0,-1.65)--node[above] {$d$}(2.4,-1.65);
  \draw[->] (1.20,0) arc[start angle=0,end angle=39.81,radius=1.20];
  \node at (1.45,.80) {$\beta$};
  \draw[->] (.70,0) arc[start angle=0,end angle=15.15,radius=.70];
  \node at (.78,-.18) {$-\alpha>0$};
  \node[align=center] at (1.15,-2.35) {The foot lies beyond the edge.\\$\alpha<0<\beta$};
\end{scope}
\end{tikzpicture}
\caption{Two possible Voronoi sectors, with the edge normal as angular
axis. In the right-hand sector both endpoints are above that axis,
so the lower endpoint parameter $\alpha$ is negative even though the
sector angle $\alpha+\beta$ is positive.}
\label{fig:signed-sectors}
\end{figure}


\subsection{Compatible flat-torus data and trace matching}

Let \(\Gamma\subset\R^2\) be a rank-two lattice and put
\(\mathbb T=\R^2/\Gamma\). Let \(N\geq1\) be an integer, let
\(v_1,\ldots,v_N\) be distinct points of \(\mathbb T\), and let
\(P\) be their full \(\Gamma\)-periodic lift. Suppose that the Voronoi
cells of \(P\) are bounded convex polygons with finitely many
face--edge incidences modulo \(\Gamma\). Assume that each incidence
has the sector form \eqref{geo:sector} and that, for one named constant
\(d_0>0\), every sector has \(d\geq d_0\). We call these
\emph{compatible periodic Voronoi data with lower edge distance \(d_0\)}.
This is conditional geometric data, not a separation assertion about
an arbitrary competitor.

The sectors partition the \(N\) cells modulo \(\Gamma\), up to their
boundary rays. Their total area is \(|\mathbb T|\), and their total
angle is \(2\pi N\): each convex cell has its site in its interior
because \(d_0>0\), so its sectors have total angle \(2\pi\).
These facts do not require an incidence bound.

Call a radius \(\eta>0\) admissible if it is smaller than the
injectivity radius of \(\mathbb T\) and, when \(N\geq2\), smaller
than half the minimum torus distance between distinct sites. Thus the
torus disks \(B(v_i,\eta)\) are embedded and pairwise disjoint; for
\(N=1\) there is no pairwise-distance condition. Write
\[
 \Omega_\eta=\mathbb T\setminus\bigcup_{i=1}^N B(v_i,\eta).
\]
In puncture limits we may further take \(\eta<d_0\), so every deleted
disk lies in its Voronoi cell. This restriction does not change a limit
as \(\eta\downarrow0\).

Fix compatible periodic Voronoi data with lower edge distance \(d_0\).
Let \(I\) be a finite or countable index set,
and choose integers \(k_i\geq6\), \(i\in I\); different indices
may have the same integer.  Suppose that
\[
 B\in C^1((0,\infty)),\qquad D_i\in C^1((0,\infty))
\]
satisfy the following assumptions:
\begin{equation}
 B(r)=-\log r+O(r^2),\qquad
 B'(r)=-1/r+O(r)
                    \quad(r\downarrow0),
 \label{dual:radial-origin}
\end{equation}
and, for every \(T_0>d_0\),
\begin{equation}
 \sum_{i\in I}\ \sup_{d_0\leq t\leq T_0}
 \bigl((1+k_i)|D_i(t)|+t|D_i'(t)|\bigr)<\infty.
 \label{dual:radial-sum}
\end{equation}
These are the exact inputs needed in the gluing and duality
argument.  The later radial construction will verify them
directly.

For a sector \eqref{geo:sector}, define \(t(\varphi)=d\sec\varphi\)
and
\begin{equation}
 H(r,\varphi)=B(r)+
       \sum_{i\in I}\left(\frac r{t(\varphi)}\right)^{k_i}
                                      D_i(t(\varphi)).
 \label{dual:trial}
\end{equation}
The same radial functions are used in every sector.

\begin{lemma}[Gluing and singular normalization]
\label{dual:gluing}
For compatible periodic Voronoi data on \(\mathbb T\) with \(N\geq1\)
distinct sites and lower edge distance \(d_0\), let \(I\) be finite or
countable, let \(k_i\geq6\) be integers with repetitions permitted,
and let \(B,D_i\in C^1((0,\infty))\) satisfy
\eqref{dual:radial-origin}--\eqref{dual:radial-sum}. The sector
functions \eqref{dual:trial} define a \(\Gamma\)-periodic
function \(H\) that is continuous off the sites and belongs to
\(H^1(\Omega_\eta)\) for every admissible \(\eta<d_0\). At every
site, uniformly in the angular variable,
\begin{equation}
 H=-\log r+O(r^2),\qquad
 \nabla H=-\frac{e_r}{r}+O(r).
 \label{dual:singularity}
\end{equation}
The gradient statement is sectorwise, and hence holds for the weak
gradient almost everywhere; no agreement of derivatives across sectors
is asserted.
In particular the constant term in \(H+\log r\) is zero.
\end{lemma}

\begin{proof}
Fix the finite collection of lifted polygons needed for one
torus. Their vertex radii are bounded by some \(T_0>d_0\). On each
closed sector, \(d_0\leq t\leq T_0\), \(r/t\leq1\), and
\(\tan\varphi\) is bounded.  The derivatives of a summand are
\[
 \partial_r\left[(r/t)^{k_i}D_i(t)\right]
       =\frac{k_i}{r}(r/t)^{k_i}D_i(t),
\]
and
\[
 \partial_\varphi\left[(r/t)^{k_i}D_i(t)\right]
       =\tan\varphi\,(r/t)^{k_i}
                         \bigl(tD_i'(t)-k_iD_i(t)\bigr).
\]
Assumption \eqref{dual:radial-sum} therefore gives uniform
convergence of the series and its first derivatives on each
closed sector away from its site.  Derivatives here are
sectorwise; derivatives across different sectors need not agree.

On a polygon edge, \(r=t\), so the trace is
\[
 B(r)+\sum_iD_i(r).
\]
The two adjacent sites have the same distance to that edge
point, and hence give the same trace.  On a ray from a site
to a vertex at distance \(r_v\), either neighboring sector has
\(t=r_v\) at its endpoint angle.  Both traces along the entire
ray are
\[
 B(r)+\sum_i(r/r_v)^{k_i}D_i(r_v).
\]
At a multiple vertex all incident sites have the same vertex
distance, so the edge traces also coincide there.  The recipe
is unchanged by period translations.

Thus all traces match.  Applying the Sobolev gluing property
to the finitely many sectors gives the asserted \(H^1\)
regularity off the holes.  No matching of normal derivatives
is required.

For \(r\leq d_0/2\), the fact that every \(k_i\geq6\), together
with \eqref{dual:radial-sum}, bounds the correction series by
\(O(r^6)\) and its gradient by \(O(r^5)\).  Combining these
bounds with \eqref{dual:radial-origin} proves
\eqref{dual:singularity}.
\end{proof}

The common notation \(D_i=b_i-u_i\), \(B=\sum_i u_i\) is one way
to supply these data.  In that notation the common edge trace
is simply \(\sum_i b_i(r)\).  The gluing argument itself only
uses \(B\) and \(D_i\).

\subsection{The exact dual identity}

Assume now that \(|\mathbb T|=N\). Let \(h\) be the mean-zero
distributional solution on \(\mathbb T\) of
\begin{equation}\label{dual:torus-potential}
 -\Delta h=2\pi\left(\sum_{i=1}^N\delta_{v_i}-1\right).
\end{equation}
The area condition is exactly the compatibility condition for this
equation. With the convention
\[
 \Gamma^*=\{k\in\R^2:k\cdot\gamma\in\Z
                     \text{ for every }\gamma\in\Gamma\},
\]
the Fourier mode \(e^{2\pi i k\cdot x}\) has eigenvalue
\(4\pi^2|k|^2\) under \(-\Delta\). The constant Fourier coefficient
of the right-hand side vanishes. Dividing its coefficient
at each nonzero \(k\in\Gamma^*\) by \(4\pi^2|k|^2\), and setting the
coefficient at zero to zero, gives a distributional solution. It is
unique because the kernel of the Laplacian on the torus consists of
constants. Define
\begin{equation}\label{dual:torus-energy}
 \mathcal W_{\mathbb T}(h)=\lim_{\eta\downarrow0}
 \left[
 \frac12\int_{\Omega_\eta}|\nabla h|^2\,\dd x+\pi N\log\eta
 \right],
\end{equation}
where the limit is taken through admissible radii. The local calculation
in Lemma~\ref{red:local}, applied in a torus chart, gives the finite
puncture limit. When \(\mathbb T=T_n\) and \(N=n\), this is exactly
\(\mathcal W_n(h)\) in \eqref{eq:torus-energy}.

For a sector \(S=S(d,\alpha,\beta)\), define
\begin{equation}
 \int_S^{\mathrm{ren}}|\nabla H|^2
 =\lim_{\eta\downarrow0}\left[
   \int_{S\cap\{r>\eta\}}|\nabla H|^2\,\dd x
                         +\theta(S)\log\eta\right].
 \label{dual:sector-ren}
\end{equation}
Lemma~\ref{dual:gluing} shows that this limit is finite.
Its counterterm uses the positive angle \(\theta(S)\), also
for a sector with a negative endpoint parameter.  Define
\begin{equation}
 J(S)=-2\pi\int_S H\,\dd x
             -\frac12\int_S^{\mathrm{ren}}|\nabla H|^2.
 \label{dual:sector-functional}
\end{equation}
The sectorwise estimates in the proof of Lemma~\ref{dual:gluing} also
show that these expressions are well defined on any individual sector
with \(d\geq d_0\) and the endpoint conditions in \eqref{geo:sector},
without first choosing a periodic tiling.

\begin{proposition}[Dual identity]
\label{dual:identity}
For compatible periodic Voronoi data on
\(\mathbb T=\R^2/\Gamma\) with \(N\geq1\) distinct sites, lower edge
distance \(d_0>0\), and \(|\mathbb T|=N\), let \(h\) be the mean-zero
potential \eqref{dual:torus-potential}. Let \(I\) be finite or countable,
let \(k_i\geq6\) be integers with repetitions permitted, and let
\(B,D_i\in C^1((0,\infty))\) satisfy
\eqref{dual:radial-origin}--\eqref{dual:radial-sum}. Let \(H\) be the
function supplied by Lemma~\ref{dual:gluing}. Then
\begin{equation}
 \mathcal W_{\mathbb T}(h)
   =\sum_S J(S)+
              \frac12\int_{\mathbb T}|\nabla(h-H)|^2\,\dd x
   \geq\sum_S J(S).
 \label{dual:energy-bound}
\end{equation}
The sum runs over all face--edge incidences in one torus.
No incidence bound or square-torus separation hypothesis is imposed.
\end{proposition}

\begin{proof}
Apply Lemma~\ref{red:local} in a local chart to \(-\nabla h\).
Near each site, \(h=-\log r+a_i+O(r)\) and
\(\nabla h=-e_r/r+O(1)\). Thus both \(h\) and \(H\) have singular
gradient \(-e_r/r\).
Their difference has square-integrable gradient.  Also
\eqref{dual:singularity} implies that
\[
 I_{\mathrm{ren}}(H)=
 \lim_{\eta\downarrow0}\left[
       \int_{\Omega_\eta}|\nabla H|^2+2\pi N\log\eta\right]
\]
exists. For admissible \(\eta<d_0\), the punctured sectors partition
\(\Omega_\eta\) up to their boundaries. The sector angles sum to
\(2\pi N\), so \(I_{\mathrm{ren}}(H)\) equals the sum of
\eqref{dual:sector-ren}.

Green's identity, with \(h\) smooth on the punctured torus and
\(H\in H^1(\Omega_\eta)\), gives
\[
 \int_{\Omega_\eta}\nabla h\cdot\nabla H
  =-2\pi\int_{\Omega_\eta}H
     +\sum_i\int_{\partial B(v_i,\eta)}
                          H\,\partial_{\boldsymbol n}h.
\]
The outward hole normal gives
\(\partial_{\boldsymbol n}h=1/\eta+O(1)\).
The zero constant term in \eqref{dual:singularity} therefore
gives
\[
 \int_{\Omega_\eta}\nabla h\cdot\nabla H
      =-2\pi\int_{\mathbb T}H-2\pi N\log\eta+o(1).
\]
Expanding \(|\nabla(h-H)|^2\), inserting this identity, and
taking the puncture limit proves
\[
 \mathcal W_{\mathbb T}(h)
   =-2\pi\int_{\mathbb T}H-\frac12I_{\mathrm{ren}}(H)
          +\frac12\int_{\mathbb T}|\nabla(h-H)|^2.
\]
Partitioning the first two terms into sectors proves
\eqref{dual:energy-bound}.
\end{proof}

If the gradient hypotheses were unchanged but the value expansion were
instead \(H=-\log r+c_i+o(1)\) at \(v_i\), the right-hand side of the
last identity would contain the additional term
\(2\pi\sum_i c_i\).  This explains why the zero constant term
in the radial construction is part of the normalization.

For a square-torus minimizer, Lemma~\ref{geo:voronoi} supplies the
compatible data with \(\Gamma=\sqrt n\,\Z^2\), \(N=n\), and
\(d_0=d_*=r_*/2=1/(2\sqrt\pi)\). Its torus has area \(n\), so
Proposition~\ref{dual:identity} applies with
\(\mathcal W_{\mathbb T}=\mathcal W_n\). The triangular application
in the next section supplies its own period and sectors and does not
use a square-torus minimizer.

\subsection{From scalar estimates to configuration bounds}

The two scalar bounds below control ordinary and signed sectors by
different arguments. Once they yield an affine lower bound for each
sector, the dual identity and the incidence count give the energy
bound. We formulate this implication for compatible flat tori first;
separation of a square-torus minimizer then removes the geometric
hypotheses from the finite-volume comparison.

\begin{lemma}[Signed endpoint parameters]
\label{dual:signed-endpoints}
Let \(K<0\) and \(d_0>0\). For each \(d\geq d_0\), let
\(G(d,\cdot)\) be an even locally integrable real function on \(\R\),
and put
\[
 g_d(x)=\int_0^xG(d,s)\,\dd s.
\]
Suppose that, for every \(d\geq d_0\),
\begin{equation}
 2g_d(x)\geq K\quad(x\geq0),\qquad
 \int_0^\infty\max\{0,-G(d,s)\}\,\dd s\leq-K.
 \label{dual:scalar-input}
\end{equation}
Then, for every \(d\geq d_0\) and every pair \(\alpha,\beta\) with
\(|\alpha|,|\beta|<\pi/2\) and \(\alpha+\beta>0\),
\begin{equation}
 g_d(d\tan\alpha)+g_d(d\tan\beta)\geq K.
 \label{dual:signed-bound}
\end{equation}
\end{lemma}

\begin{proof}
Evenness of \(G(d,\cdot)\) makes \(g_d\) odd. If
\(\alpha,\beta\geq0\), the first inequality in
\eqref{dual:scalar-input} bounds each term in
\eqref{dual:signed-bound} below by \(K/2\).

If \(\alpha<0\), set
\[
 u=-d\tan\alpha,\qquad v=d\tan\beta.
\]
The inequality \(\alpha+\beta>0\) gives \(\beta>-\alpha>0\).
Since the tangent is increasing on \((-\pi/2,\pi/2)\), we have
\(0<u<v\). Oddness and the second inequality in
\eqref{dual:scalar-input} give
\[
 \begin{aligned}
 g_d(d\tan\alpha)+g_d(d\tan\beta)
 &=g_d(v)-g_d(u)=\int_u^vG(d,s)\,\dd s\\
 &\geq-\int_0^\infty\max\{0,-G(d,s)\}\,\dd s\geq K.
 \end{aligned}
\]
The case \(\beta<0\) is symmetric. Both endpoints cannot be negative
because their sum is positive, so these cases are exhaustive.
\end{proof}

\begin{proposition}[From sector bounds to energy bounds]
\label{dual:assembly}
Let \(d_0>0\), let \(I\) be finite or countable, let \(k_i\geq6\)
be integers with repetitions permitted, and let
\(B,D_i\in C^1((0,\infty))\) satisfy
\eqref{dual:radial-origin} and \eqref{dual:radial-sum} with
this lower distance \(d_0\). For every sector \eqref{geo:sector} with
\(d\geq d_0\), use these same radial data in \eqref{dual:trial} and
let \(J(S)\) be the functional \eqref{dual:sector-functional}.
Suppose that \(K\leq0\) and \(\lambda,\mu\in\R\) satisfy
\begin{equation}
 \begin{aligned}
 J(S(d,\alpha,\beta))
 &\geq K+\lambda A(S(d,\alpha,\beta))\\
 &\qquad+\mu\theta(S(d,\alpha,\beta)),
 \end{aligned}
 \label{dual:sector-input}
\end{equation}
for every \(d\geq d_0\), \(|\alpha|,|\beta|<\pi/2\), and
\(\alpha+\beta>0\).

\emph{Compatible tori.}
Let \(\mathbb T=\R^2/\Gamma\) have area \(N\), where \(N\geq1\)
is an integer. Suppose \(N\) distinct sites supply compatible periodic
Voronoi data with lower edge distance \(d_0\) and at most \(6N\)
face--edge sectors. Their mean-zero potential
\eqref{dual:torus-potential} satisfies
\begin{equation}\label{dual:compatible-bound}
 \mathcal W_{\mathbb T}(h)\geq N(6K+\lambda+2\pi\mu).
\end{equation}

\emph{All square-torus configurations.}
Assume in addition \(d_0\leq d_*=1/(2\sqrt\pi)\). For every
integer \(n\geq2\) and every configuration of \(n\) distinct
points \(v_1,\ldots,v_n\) on \(T_n=\R^2/(\sqrt n\,\Z^2)\), let \(h\)
be the mean-zero solution of
\[
 -\Delta h=2\pi\left(\sum_{i=1}^n\delta_{v_i}-1\right).
\]
Then its puncture-first energy \eqref{eq:torus-energy}, with embedded
and pairwise disjoint punctures, satisfies
\begin{equation}
 \mathcal W_n(h)\geq n(6K+\lambda+2\pi\mu).
 \label{dual:all-configurations}
\end{equation}
More generally, if a periodic locally integrable field \(E\) on \(T_n\)
satisfies, distributionally,
\[
 \operatorname{div}E=2\pi\left(\sum_{i=1}^n\delta_{v_i}-1\right),
 \qquad \operatorname{curl}E=0,
\]
then \(E=-\nabla h+c\) for a constant \(c\in\R^2\), and its energy
with the same punctures and counterterm obeys
\[
 \mathcal E_{T_n}(E)=\mathcal W_n(h)+\frac n2|c|^2
                    \geq n(6K+\lambda+2\pi\mu).
\]
The energy identity remains valid after applying the same fixed
orthogonal rotation to both fields.
\end{proposition}

\begin{proof}
For compatible data, the sectors have total area \(N\) and total
angle \(2\pi N\). Write \(M_0\leq6N\) for their number.
Proposition~\ref{dual:identity} and \eqref{dual:sector-input} give
\[
 \begin{aligned}
 \mathcal W_{\mathbb T}(h)
 &\geq\sum_S J(S)\\
 &\geq M_0K+\lambda N+2\pi\mu N\\
 &\geq N(6K+\lambda+2\pi\mu).
 \end{aligned}
\]
The last inequality uses \(M_0\leq6N\) and \(K\leq0\); it includes
\(K=0\), when both terms involving \(K\) vanish.

Now assume \(d_0\leq d_*\) and fix \(n\geq2\).
Lemma~\ref{geo:existence} supplies a minimizing configuration of
distinct points, with potential \(h_{\min}\).
Lemma~\ref{geo:separation} separates its full periodic lift by
\(r_*=1/\sqrt\pi\). Lemma~\ref{geo:voronoi} then gives at most
\(6n\) sectors, each with edge distance at least \(d_*\geq d_0\).
These are compatible data on \(T_n\), so
\eqref{dual:compatible-bound} applies to \(h_{\min}\).
Every other distinct configuration has energy at least
\(\mathcal W_n(h_{\min})\), proving
\eqref{dual:all-configurations} without a separation assumption on
that configuration.

Equation~\eqref{red:harmonic} applies to the displayed divergence and
curl conditions and gives the constant harmonic component and its
exact energy penalty. A fixed orthogonal rotation preserves the
squared norm and the puncture counterterm, giving the last assertion.
\end{proof}

The next section chooses particular radial data and exact constants
\(K<0,\lambda,\mu\). For these data, the primitive identity
\eqref{cal:sector-primitive} and the two scalar bounds yield
\[
 J(S)\geq K+\lambda A(S)+\mu\theta(S)
\]
for every relevant sector. Proposition~\ref{cal:scalar-to-torus}
applies the compatible-torus conclusion to this construction and
identifies \(6K+\lambda+2\pi\mu\) with the triangular energy.
Section~\ref{sec:conclusion} uses the square-torus conclusion, whose
minimizer argument has already removed the separation hypothesis.
